\section{Γενετικοί Αλγόριθμοι και Μοντέλο Νήσων}
\label{sec:bg_ga_islands}

Οι Γενετικοί Αλγόριθμοι (GAs) ανήκουν στην κατηγορία των μεθευρετικών μεθόδων καθολικής βελτιστοποίησης, εμπνευσμένοι από τη δαρβινική φυσική επιλογή~\cite{goldberg1989genetic}. Τυπικά, ένα \textit{χρωμόσωμα} (chromosome) αναπαριστά μία υποψήφια λύση ως διάνυσμα πραγματικών τιμών $\bm{x} = (x_1, x_2, \ldots, x_n) \in \mathbb{R}^n$, ενώ η \textit{συνάρτηση καταλληλότητας} (fitness function) $f : \mathbb{R}^n \to \mathbb{R}$ αποδίδει σε κάθε άτομο μία ποσοτική αξιολόγηση της ποιότητάς του. Στόχος του αλγορίθμου είναι η εύρεση $\bm{x}^* = \arg\min f(\bm{x})$ χωρίς να απαιτείται η πληροφορία της κλίσης $\nabla f$.

\paragraph{Επιλογή με τουρνουά.}
Σε κάθε γενιά επιλέγονται γονείς μέσω τουρνουά μεγέθους $k$: από $k$ τυχαία επιλεγμένα άτομα του πληθυσμού, επιβιώνει εκείνο με την υψηλότερη καταλληλότητα. Η πιθανότητα επιλογής ενός ατόμου με κατάταξη $r$ (από τον χειρότερο προς τον καλύτερο) εκτιμάται ως:
\begin{equation}
    P_{\text{sel}}(r) \approx \frac{k}{N}\binom{N-1}{k-1}^{-1}\binom{N-r}{k-1},
\end{equation}
όπου $N$ είναι το μέγεθος πληθυσμού. Για $k=2$ (δυαδικό τουρνουά), η επιλεκτική πίεση είναι ήπια, αποτρέποντας την πρόωρη σύγκλιση.

\paragraph{Διασταύρωση SBX.}
Η Simulated Binary Crossover (SBX) μιμείται τη μονοσημειακή διασταύρωση δυαδικής κωδικοποίησης στον χώρο των πραγματικών αριθμών~\cite{goldberg1989genetic}. Από δύο γονείς $x_i^{(1)}, x_i^{(2)}$, τα δύο απόγονα υπολογίζονται ως:
\begin{equation}
    x_i^{(1,2)'} = \tfrac{1}{2}\bigl[(1 \pm \beta_i)(x_i^{(1)} + x_i^{(2)})\bigr],
\end{equation}
όπου ο παράγοντας εξάπλωσης $\beta_i$ ορίζεται μέσω παραμέτρου κατανομής $\eta_c \geq 0$: για μικρές τιμές $\eta_c$ παράγονται απόγονοι μακριά από τους γονείς (εξερεύνηση), ενώ για μεγάλες τιμές οι απόγονοι παραμένουν κοντά (εκμετάλλευση).

\paragraph{Μετάλλαξη Gauss.}
Κάθε γονίδιο $x_i$ μεταλλάσσεται με πιθανότητα $p_m = 1/n$ προσθέτοντας διαταραχή Gauss:
\begin{equation}
    x_i' = x_i + \mathcal{N}(0,\, \sigma_m^2),
\end{equation}
όπου $\sigma_m$ ελέγχει το πλάτος αναζήτησης. Το αποτέλεσμα αποκόπτεται στα όρια $[x_i^{\min}, x_i^{\max}]$ ώστε να παραμείνει εντός του εφικτού χώρου.

\paragraph{Ελιτισμός $(\mu + \lambda)$.}
Στη στρατηγική $(\mu + \lambda)$, οι $\mu$ γονείς και τα $\lambda$ νέα απόγονα αντιμετωπίζονται μαζί, και τα $\mu$ καλύτερα άτομα επιλέγονται για την επόμενη γενιά. Έτσι εγγυάται η μονοτονική βελτίωση της βέλτιστης λύσης ανά γενιά, αφού κανένα ήδη ανακαλυφθέν καλό άτομο δεν απορρίπτεται.

\paragraph{Μοντέλο Νήσων (Island Model DGA).}
Στο μοντέλο νήσων~\cite{whitley1999island}, ο συνολικός πληθυσμός $N_{\text{total}}$ διαιρείται σε $K$ υποπληθυσμούς μεγέθους $N_s = N_{\text{total}} / K$, οι οποίοι εξελίσσονται ανεξάρτητα σε ξεχωριστούς υπολογιστικούς κόμβους. Σύμφωνα με την ανάλυση του Cantú-Paz~\cite{cantupaz2000parallel}, το βέλτιστο μέγεθος υποπληθυσμού εξαρτάται από τη διακύμανση της καταλληλότητας και τον ρυθμό μετανάστευσης: μεγάλοι υποπληθυσμοί συγκλίνουν ανεξάρτητα, ενώ μικροί υποπληθυσμοί χρειάζονται συχνότερη ανταλλαγή ατόμων για να αποφύγουν τη γενετική απώλεια (genetic drift). Οι Alba \& Troya~\cite{alba1999survey} κατέγραψαν εμπειρικά ότι $K \in [4, 16]$ νήσοι με ρυθμό μετανάστευσης $m \approx 5{-}10\%$ ανά $M \in [10, 30]$ γενιές επιτυγχάνουν ευρωστία και κλιμακωσιμότητα.

Περιοδικά, ανά $M$ γενιές, επιλεγμένα άτομα (μετανάστες) αποστέλλονται σε γειτονικές νήσους σύμφωνα με μια \textit{τοπολογία δακτυλίου}:
\begin{equation}
    \mathcal{P}_{i}^{(t+M)} \leftarrow \text{SelectElites}\!\left(\mathcal{P}_{i}^{(t+M)}\right) \cup \text{ReceiveMigrants}\!\left(\mathcal{P}_{i-1}^{(t+M)}\right).
\end{equation}
Η τοπολογία δακτυλίου επιλέγεται γιατί, όπως αποδείχθηκε από τους Doerr, Friedrich et al.~\cite{friedrich2017island}, η διάδοση πληροφορίας (rumor spreading) σε κυκλικό γράφο επιτρέπει βέλτιστη εξισορρόπηση μεταξύ εξερεύνησης και εκμετάλλευσης: η πληροφορία ενός καλού ατόμου διαδίδεται σε $O(\log K)$ γενιές σε όλες τις νήσους, αποτρέποντας την ταυτόχρονη σύγκλιση σε διαφορετικά τοπικά ακρότατα.

\paragraph{Ανίχνευση στασιμότητας.}
Για την αποφυγή σπατάλης υπολογιστικών πόρων, ο αλγόριθμος παρακολουθεί την τυπική απόκλιση $\sigma_f$ της καταλληλότητας στον πληθυσμό. Εάν $\sigma_f < \epsilon_{\text{stag}}$ για $G_{\text{stag}}$ συνεχόμενες γενιές, ο υποπληθυσμός θεωρείται στάσιμος (stagnated) και εκκινείται επανεκκίνηση (restart) με τυχαία νέα άτομα, διατηρώντας μόνο την κορυφαία ελίτ.
